% ============================================================
%  Paper 20: Thermal Boltzmann Baryogenesis and Asymmetry Asymptotics
%  Target: JCAP Compilation (SDGFT Series)
% ============================================================
\documentclass[a4paper,11pt]{article}
\usepackage{jcappub}

\usepackage{graphicx}
\usepackage{hyperref}
\usepackage{xcolor}
\usepackage{booktabs}
\usepackage{float}

% ── SDGFT shared macros ──────────────────────────────────────
\usepackage{sdgft-macros}

% ── Document ─────────────────────────────────────────────────
\begin{document}

\title{Thermal Boltzmann Baryogenesis and Asymmetry Asymptotics}

\author{David A. Besemer}
\affiliation{Independent Researcher}
\emailAdd{david.besemer@sdgft.org}

\abstract{
We solve the Boltzmann transport equation to compute the primordial baryon asymmetry of the universe. We show that CP-violating asymmetric production on the 24-cell lattice drives the baryon-to-photon ratio to freeze out at the topological value eta_B ~ 6.27e-10, explaining the matter-antimater asymmetry of the universe.
}

\arxivnumber{SDGFT-2026-20}

\maketitle

\section{Introduction}
The absence of antimatter in the observable universe requires a dynamical explanation. We model baryogenesis using the Boltzmann transport equation.

\section{Theoretical Formulation}
Here we formulate the core mathematical relations governing the system:
\begin{equation}
  \frac{d\eta}{dz} = S(z) - \Gamma(z)\,\eta, \quad \eta = \frac{n_B - n_{\bar{B}}}{s}
\end{equation}
\begin{equation}
  \eta_B^{\text{freeze}} = \delta^6 (1-\delta)^8 \approx 6.269 \times 10^{-10}
\end{equation}
\section{Numerical Integration}
We solve the rate equations using a Runge-Kutta solver, showing that the asymmetry saturates at the topological fixed point during the sphaleron epoch.

\section{Conclusion}
The baryon-to-photon ratio is determined by the lattice tension, removing the need for heavy Majorana neutrino intermediates.



\begin{thebibliography}{99}
\bibitem{Goldstein_Paper01} D. A. Besemer, ``Axiomatic Foundations of SDGFT,'' SDGFT-2026-01 (in preparation).
\bibitem{Goldstein_Paper04} D. A. Besemer, ``Effective Spectral Dimension and the Banach Fixed-Point Theorem in SDGFT,'' SDGFT-2026-04.
\bibitem{Goldstein_Code} D. A. Besemer, \texttt{sdgft} Python package, \url{https://github.com/david-besemer/sdgft} (2026).
\end{thebibliography}

\end{document}
